\documentclass{article}
\usepackage{amsfonts}
\usepackage{amsmath}
\usepackage{amssymb}
\usepackage{mathrsfs}
\usepackage{xcolor}
\usepackage{graphicx}
\newtheorem{theorem}{Theorem}
\newtheorem{lemma}[theorem]{Lemma}
\newtheorem{proposition}[theorem]{Proposition}
\newtheorem{definition}[theorem]{Definition}

\def\det{\mathrm{det}}
\def\nondet{\mathrm{rnd}}
\newcommand{\term}{\textbf}

\title{Evolutionary Fitness and the Identical Ancestor Point}

\author{Samuel Allen Alexander}

\begin{document}


\maketitle

\begin{abstract}
    ...
\end{abstract}

\section{Introduction}

The idea that evolutionary fitness is synonymous with having more children
is almost a cliche because it is invoked so often.
But this naive approach does not do full justice to the nuances of fitness.
If I give birth to a million sterile children (and no fertile children),
then that should not make me fit whatsoever. Giving birth to one fertile child
should confer more fitness than giving birth to any number of sterile children.

We could attempt to patch the above problem by considering grandchildren
rather than children: measure evolutionary fitness by how many grandchildren
an organism has. But this just pushes the bubble around under the wallpaper,
because one could theoretically have numerous sterile grandchildren (and no
fertile grandchildren), which would also make one manifestly unfit. Delaying
the calculation any number of generations runs into the same problem.

To be clear, the problem does not merely arise in situations like ``all of this
organism's children are sterile,'' which may seem contrived.
The problem can arise in very realistic scenarios. For example, as noted by ...,
organisms could inhabit an environment which very rarely suffers from heatwaves.
If the heatwaves are rare enough, organisms could evolve to lose expensive heatwave
resistance, which would seem (in the short term) to make them more fit.
But then if all such mutants are killed by the next heatwave, then the
mutation was actually unfit.

Another example: naive ``fitness $=$ number of children'' examples often gloss over
the fact that in nature, organisms can sexually reproduce with their own children.
It is not clear whether having a child with its own child should grant the same
amount of fitness to a fruitfly as having a child with an unrelated fruitfly mate.

We address this by using the idea of identical ancestor points from
... and ..., more recently studied by .... Motivated by the idea of identical
ancestor points, we consider two types of organisms. We call an organism $x$ a
\emph{winner} if, at some point in time (probably in the distant future after
$x$'s death), all living members of $x$'s species are descended from $x$.
We call $x$ a \emph{loser} if $x$'s descendants eventually go extinct.
We argued ... for what we call the \emph{identical ancestor point axiom}.
This axiom states that in any species, every organism of that species is either
a winner or a loser, in the above sense. This offers a shockingly simple
new way to measure the fitness of a finite set $S$ of organisms in a
species, assuming the species satisfies the axiom. Namely: the fitness of
$S$ is
\[
    \frac{\mbox{number of winners in $S$}}{|S|}.
\]

In this paper we discuss empirical results about such fitness based on
simulations \emph{in silica}.

The roadmap of this paper is as follows:

\begin{itemize}
    \item
    In Section ..., we discuss the theory behind this proposed new methodology
    in greater detail.
    \item
    In Section ..., we describe a bespoke simulation which we implemented
    in Python in order to empirically measure evolutionary fitness values.
    \item
    In Section ..., we describe empirical results obtained from the aforementioned
    simulation.
    \item
    In Section ..., we summarize and make concluding remarks.
\end{itemize}

\section{Theory behind this new methodology}

The so-called \emph{identical ancestor point axiom} states that any species can be
partitioned into so-called \emph{winners} and \emph{losers}, defined in the
Introduction. This is not useful for studying the fitness of any particular
organism, because any particular organism is either a winner or a loser within
its species, and so would be given fitness value of either $0$ or $1$, which is
too granular to be of any use. Thus, we instead measure the fitness of sets of
organisms in a species.

Sets of organisms correspond naturally to actions, in the sense that
to any action, we can associate the set of organisms who take that action.
For example, consider the action of eating a poisonous mushroom.
This action can be associated with the set which contains precisely those
organisms who ate a poisonous mushroom during their lifetime.

Absolute values computed using this methodology are, as far as we know
currently, not meaningful by themselves (save only in degenerate cases
such as $0$ or $1$). Rather, what are meaningful are the values relative
to other values. In other words, knowing that the set of organisms who
eat poisonous mushrooms has a fitness score of $0.65$ is not particularly
meaningful in itself, but can be compared to other scores. If the fitness
score of eating delicious strawberries is $0.85$ then since
$0.85>0.65$, we conclude that eating delicious strawberries is more
evolutionarily fit than eating poisonous mushrooms.

\section{A python simulation}

This paper is intended as a proof of concept for a new methodology for measuring
evolutionary fitness. We implemented a simplisitic simulation of organisms using
python. In this simulation, a small 2-dimensional world is initially populated
with a random number of randomly scattered organisms, all of the same species.
Time in this world is discrete, broken up into \emph{turns}. On every turn,
every organism takes one \emph{action}. Possible actions are:
\begin{itemize}
    \item Attempt to walk north.
    \item Attempt to walk south.
    \item Attempt to walk east.
    \item Attempt to walk west.
    \item Attempt to eat adjacent food.
    \item Attempt to reproduce with adjacent fellow organism.
\end{itemize}
A $Q$-learner reinforcement learning (RL) agent was used to choose which actions
each agent took. Observations were passed to the RL agent as if the organisms
had severely limited eyesight, in order to facilitate quick training of the
agent ($Q$-learners are very quick to act, but the time it takes them to learn
their environment grows exponentially in the eyesight of the agents).
This means the organisms acted purely on instinct with no conscious memories,
choosing each action solely based upon what they saw at that moment and not
directly upon what they had seen earlier (though the reinforcement learning
agent did train to update its weights, so that the organisms did possess a
kind of communical instinctual memory).

Because this paper is only intended as a simple proof of concept, we will not
go into fine details here about the exact parameters used in the simulation;
the interested reader can consult the code itself, available on Github ....
One particular detail which is important for understanding the results we
obtained concerns how reproduction works in this simulation. In our simulation,
when an organism attempts to reproduce with an adjacent fellow organism, the
attempt will fail unless both of the following conditions are true:
\begin{enumerate}
    \item There is \emph{exactly one} other organism adjacent to the organism
        who is attempting to reproduce.
    \item There is at least one empty space adjacent to the organism
        who is attempting to reproduce (so that the child, which is born
        immediately, has a place to be born into).
\end{enumerate}
We implemented the simulation this way in order to simplify many troublesome
questions regarding parenthood and so on, while keeping the available action
list manageable. As a side effect, this causes crowding to naturally reduce
fertility, since, in a large crowd, it is difficult to satisfy condition 1.
We will see in the next section that this apparently resulted in a non-increasing
child-count vs.\ fitness score: beyond a certain point, having more children
was measured to actually confer \emph{less} fitness, probably due to this
crowding mechanism. While this is of course contrived and unnatural, it
nevertheless points to similar phenomena plausible in real life: beyond a certain
point, having too many children really could be detrimental to one's fitness,
if those children are competing for limited resources, etc.

Within the simulation, food periodically grows at random locations in the
environment. This food is split into eight different types (below, Strl should
be read as ``sterile''; Fut should be read as ``future'';
Desc should be read as ``descendant'';
Frt should be read as ``fertility bonus''):
\begin{itemize}
    \item $\mathsf{Normal}$ food has no effect other than default satiation.
    \item $\mathsf{StrlFutDesc}$ food is a food which has the
        following effect. After eating the food, any time that the eater
        subsequently reproduces, the offspring will be born sterile.
        Said offspring will never be succeed at reproducing. Thus, if one eats
        this food before having
        any offspring, then one cannot ever have any grandchildren.
    \item $\mathsf{StrlAllDesc}$ food is a food which causes
        all of the eater's descendants to become sterile, including
        descendants born after the food is eaten, but also including
        descendants born \emph{before} the food was eaten (of course
        this is physically unrealistic, but we included it because
        it illuminates how our fitness measure works).
    \item $\mathsf{StrlFutDescPlusFrt}$ food is exactly the same
        as $\mathsf{StrlFutDesc}$ food, except that it additionally
        confers an added fertility bonus to the eater, so that the eater
        will enjoy easier reproduction after eating it. (However, the offspring
        of this easier reproduction will be sterile.)
    \item $\mathsf{StrlAllDescFrt}$ food is exactly the same
        as $\mathsf{StrlAllDesc}$ food, except it additionally confers an
        added fertility bonus, as in the previous bullet.
    \item $\mathsf{Frt}$ food is a food which grants additional fertility to
        its eater (in exactly the same way as in the previous two bullets).
        It has no other side effects.
    \item $\mathsf{Insanity}$ food is a food which hijacks the
        reinforcement learning
        agent driving the eater. After eating insanity food, the eater will
        thereafter act in ways which maximize punishment and minimize reward
        (rather than the usual strategy of minimizing punishment and
        maximizing reward). This is partly inspired by ....
    \item $\mathsf{Shortsightedness}$ food is a food which alters the observations
        which are sent to the reinforcement learning agent in order to choose
        the eater's actions. To make a long story short, this food has the
        effect of reducing the eater's range of vision.
\end{itemize}

\section{Empirical results from the simulation}

We ran the simulation 2488 times, terminating each run after 100,000 turns
if the species did not go extinct before then. Of the runs where the species
went extinct, extinction always occurred very early: if the species survived
more than 200 turns, then it never was observed to go extinct before 100,000
turns.

For each non-extinct run, we tracked which organisms were winners and which
organisms were losers by the time we terminated the run. For some organisms born
late in the run, there was not enough time to determine their winner/loser
status (they had living descendants and living non-descendants at the time
when the run was terminated). These organisms were discarded for purposes of
statistical analysis. From the organisms not discarded, we calculated fitness
scores for each food type, and also for each childcount. The food type fitness
scores we calculated are as follows:

\begin{table}[ht]
    \centering
    \caption{Evolutionary fitness scores by food type.}
    \label{tab:fitness-by-food}
    \begin{tabular}{lr}
        \hline
        Food type & Evolutionary fitness score \\
        \hline
        $\mathsf{Normal}$ & 0.7658836489 \\
        $\mathsf{StrlFutDescs}$ & 0.4764546404 \\
        $\mathsf{StrlAllDescs}$ & 0.0000000000 \\
        $\mathsf{StrlFutDescsFrt}$ & 0.4748666463 \\
        $\mathsf{StrlAllDescsFrt}$ & 0.0000000000 \\
        $\mathsf{Frt}$ & 0.8214282319 \\
        $\mathsf{Insanity}$ & 0.6727775267 \\
        $\mathsf{Shortsighted}$ & 0.7658694215 \\
        \hline
    \end{tabular}
\end{table}

We make the following observations:
\begin{enumerate}
    \item $\mathsf{Normal}$ food is more fit than any food type with any detrimental effect.
    \item $\mathsf{FrtlBonus}$ food is the most fit of all the foods.
    \item $\mathsf{StrlAllDescs}$ and $\mathsf{StrlAllDescsBonus}$ had
        fitness scores of $0$.
    \item $\mathsf{StrlFutDescs}$ and $\mathsf{StrlFutDescsFrt}$ had
        almost identical fitness scores, the latter being slightly smaller
        than the former.
    \item $\mathsf{Insanity}$ was worse than normal food but not as bad
        as any of the sterility side-effect foods.
    \item $\mathsf{Shortsighted}$, surprisingly, had virtually the exact
        same fitness as $\mathsf{Normal}$.
\end{enumerate}

Most of these findings are unsurprising. The $0$ fitness of the
foods which confer sterility to all descendants corresponds to the fact that
eating such foods guarantees one will become a loser, regardless of how fertile
one is. This empirically confirms our observation in the Introduction that
having many sterile children (and no non-sterile children) is inherently unfit,
contrary to the naive approach of measuring fitness solely by counting
offspring.

The following findings are rather surprising:
\begin{itemize}
    \item $\mathsf{StrlFutDescs}$ actually outperforms
        $\mathsf{StrlFutDescsFtl}$. The amount by which it outperforms
        is quite small; this could just be a statistical fluke.
        But it does make sense that $\mathsf{StrlFutDescs}$ would outperform,
        because of the way reproduction works in the simulation, as remarked
        in the previous section. When one eats food which makes one's future
        descendants become sterile, and also grants one a fertility bonus
        (so that one can more prolifically produce such sterile future children),
        this is detrimental to any non-sterile children one has already
        produced before eating the food. The greater abundance of sterile
        children resulting from the fertility bonus compete with the
        fertile children born before the food was eaten.
    \item $\mathsf{Shortsighted}$ seems to have the same fitness as
        $\mathsf{Normal}$ food. We cannot see any way of explaining this,
        except merely that it turns out the longer line of sight was not
        actually beneficial in the environment in question. This illustrates
        that these fitness scores can reveal facts which are not evident
        at first glance. In this case, the fitness scores reveal that
        the superior vision of organisms that did not eat $\mathsf{Shortsighted}$
        food is not actually that helpful. Of course this would not generalize
        to arbitrary environments; it is an artifact of the particular laws
        of physics in our simulation.
\end{itemize}



\section{Summary and conclusion}



% \bibliographystyle{plain}
% \bibliography{bio}
\end{document}
